% !TEX program = pdflatex
\documentclass[a4paper,twocolumn,10pt]{article}

\usepackage[top=30mm,bottom=25mm,left=20mm,right=20mm]{geometry}
\setlength{\columnsep}{7mm}
\usepackage{CJKutf8}
\DeclareFontShape{C70}{min}{b}{n}{<->ssub*min/m/n}{}
\DeclareFontShape{C70}{min}{bx}{n}{<->ssub*min/m/n}{}
\usepackage{graphicx}
\usepackage{booktabs}
\usepackage{array}
\usepackage{url}
\usepackage{titlesec}
\usepackage{setspace}
\titlespacing*{\section}{0pt}{1.1ex plus .3ex minus .2ex}{0.6ex}
\titlespacing*{\subsection}{0pt}{0.8ex plus .2ex minus .2ex}{0.4ex}
\titleformat{\section}{\large\bfseries\sffamily}{\thesection}{0.7em}{}
\titleformat{\subsection}{\normalsize\bfseries\sffamily}{\thesubsection}{0.6em}{}
\setstretch{1.06}
\pagestyle{empty}

\newcommand{\code}[1]{{\small\ttfamily\bfseries #1}}

\begin{document}
\begin{CJK}{UTF8}{min}

\title{\vspace{-8mm}SOC調査起点からの証跡付き行動列復元に向けた\\
エージェント型ログ探索手法の評価\\[1mm]
{\large Evaluation of Agent-based Log Investigation for Evidence-backed Behavior Reconstruction from SOC Investigation Clues}}
\author{著者名$^{1}$ \quad 共著者名$^{2}$\\
{\small $^{1}$所属名 \quad $^{2}$所属名}}
\date{}
\maketitle
\thispagestyle{empty}

\begin{abstract}
Security Operation Center (SOC) analysts must connect alerts and endpoint telemetry into an evidence-backed behavior story. This paper evaluates a CLOUSEAU-based LLM agent workflow for reconstructing behavior chains from ATLASv2-derived Windows endpoint logs. We use 27 benign behavior chains and three staged input conditions: alert-assisted, process/time clue, and process/time clue without alert summary rows. The stronger model reconstructed 83.9\% of required behavior items and 80.9\% of critical evidence items across 81 runs, while the smaller model showed weaker evidence grounding and ordering. The results clarify the role of alert summaries and the remaining challenges of overclaim control and sequence construction.
\end{abstract}

\section{はじめに}

SOC（Security Operation Center）におけるインシデント分析では，検知アラートを受け取った後に，端末上で実際に何が起きたかを説明可能な形で再構成する必要がある．アラート名や検知理由は調査の入口として有用である一方，それだけでは親子プロセス，コマンドライン，ファイル・レジストリ・ネットワーク操作，それらを支えるログ行の対応関係を十分に説明できない．調査者が最終的に必要とするのは，「怪しい」というラベルではなく，観測証跡に基づいて，どの主体が，どの対象に，どの順序で操作したかを示す行動列である．

本研究が対象とするのは，マルウェア検知そのものではなく，検知後のエンドポイント調査である．特に，SOCが少数の手がかりから調査を開始する状況を想定する．入力は，CBC alertの要約，またはhost，process，time window程度のprocess/time clueであり，出力は証跡付きの行動ステップ列である．この設定では，CLOUSEAUのようなLLM階層型調査エージェント\cite{clouseau}が，自然言語の手がかりから探索方針を作り，SQLite化されたログに対して反復的に問い合わせを行うことで，関連する観測行を集められる可能性がある．しかし，アラート理由を言い換えただけの出力，近傍の無関係な行動の混入，証跡で支えられない推定，行動順序の誤りが生じる危険もある．

本稿では，ATLASv2 benign H1 / benign-1由来のWindowsホストログを用い，27個の良性セキュリティ行動chainを対象に，CLOUSEAUを基盤とするエージェント型ログ探索手法を評価する．ここで既存研究との立ち位置を明確にしておく．CLOUSEAUは，Point-of-Interest（POI）から攻撃調査を開始し，Chief Inspector，Investigator，Question-Answering agentが協調してattack narrativeを復元する階層型multi-agent frameworkである\cite{clouseau}．ATLASはaudit logから攻撃ストーリを復元する攻撃調査支援手法であり\cite{atlas}，ATLASv2はその攻撃エンゲージメントと通常業務活動を拡張したデータセットである\cite{atlasv2}．本研究はCLOUSEAUやATLASを置き換える攻撃調査アルゴリズムではない．CLOUSEAUの調査パイプラインとATLASv2のログ環境を用い，SOCの検知後調査を模した行動列復元ベンチマークとして，アラート要約への依存度，低レベルtelemetryからの復元可能性，証跡同定の成否を測る研究である．

本稿の貢献は次の3点である．第一に，CLOUSEAUのPOI起点調査をSOCの検知後行動復元へ読み替え，行動item，critical evidence，sequence order，candidate precision，overclaimに分解した評価枠組みとして定義する．第二に，27 chainを3段階の入力条件で実行し，アラート支援条件，process/time clue条件，アラート要約行を隠したtelemetry中心条件を比較する．第三に，ATLASv2由来ログを用いた検知後インシデント分析の位置づけを明確化し，明示的実行チェーン，複数段ツールチェーン，意味解釈型チェーンごとの失敗傾向を分析する．

\section{関連研究と課題設定}

\subsection{検知後調査としての行動列復元}

エンドポイント調査では，Windows Securityログ，Sysmon，EDR telemetry，DNS，ブラウザ履歴など，粒度や記録目的の異なるログを横断する．Sysmonはプロセス生成，親プロセス，コマンドライン，ネットワーク接続，レジストリ操作などを記録できる\cite{sysmon}．EDRは検知理由やプロセス情報をまとめたalert summaryを提供する場合がある．これらは調査に有用だが，alert summaryは観測された全行動を代表するものではない．したがって，最終出力はアラート文面ではなく，ログ上で確認できる行動主張に限定する必要がある．

本研究の課題は，異常検知モデルの精度評価ではない．DeepLogのような系列異常検知研究は，ログ系列から異常を検出・診断することに焦点を置く\cite{deeplog}．一方，本研究は検知後のSOC調査において，調査者が検証可能な証跡付き行動列を復元できるかを問う．この違いは評価単位にも表れる．本研究では自然文の類似度ではなく，主体，操作，対象，コマンドライン，証跡行などの必須itemを採点する．

\subsection{LLMエージェントによる探索}

LLMエージェントは，初期手がかりから仮説を立て，外部環境への問い合わせ結果に応じて次の探索を変更できる．ReActは推論と行動を交互に進める枠組みを示しており\cite{react}，ログ調査でも，時刻範囲，プロセス木，ファイル・レジストリ対象を段階的に確認する流れと親和性がある．

セキュリティ分野では，CLOUSEAUが階層型multi-agentによる自律的な攻撃調査を提案している\cite{clouseau}．CLOUSEAUは，raw incident logsをLLMが扱いやすい構造化形式へ変換し，QA agentを通じてログを自然言語的に検索可能にし，Investigator agentとChief Inspector agentが調査leadの生成，証跡収集，attack storyの統合を行う．本研究はこの設計思想に基づくが，評価対象を「攻撃story全体の復元」から「SOCアラートやprocess/time clue周辺の良性を含む行動chainの証跡付き復元」へ移す．そのため，最終判定の正否ではなく，観測ログで支えられた行動itemと過剰主張を採点する．

ただし，セキュリティログ調査では，もっともらしい説明は十分ではない．エージェントが関連しそうなログを広く集めるほど，主行動列に属さない近傍行動や，同じ操作の重複表現も出力しやすい．そのため本研究では，正しいitemを拾えたかだけでなく，余分な主張をどれだけ抑えたか，critical evidenceを対応付けられたか，行動順序を構成できたかを分けて評価する．

\subsection{ATLAS/ATLASv2の位置づけ}

ATLASは，audit logに基づいて攻撃ストーリを復元するため，因果依存グラフと系列学習を組み合わせた攻撃調査支援手法である\cite{atlas}．ATLASv2は，ATLASの攻撃シナリオを再現しつつ，通常業務活動とSysmonおよびCarbon Black Cloud由来のログを拡充したデータセットとして報告されている\cite{atlasv2}．本研究では，ATLASv2を「攻撃検知済みの答えを読むデータ」ではなく，「SOCでアラートやtelemetryを受け取った後，端末上の行動を証跡付きで説明する能力を測るデータ」として利用する．

この立ち位置により，本研究の対象インシデント分析は明確になる．対象は，侵害全体の帰属や最終判定ではなく，アラート周辺の端末行動復元である．たとえばDNS packet capture，Python SimpleHTTPServer，Sublime経由のscript execution，cmd.exe周辺操作，Discord Run Key関連操作など，良性環境内でもセキュリティアラートや調査対象になり得る行動について，観測ログから説明可能なchainを作れるかを評価する．

\section{提案手法}

\subsection{パイプライン}

図\ref{fig:pipeline}に，本研究で用いるCLOUSEAU型のエージェント型ログ探索パイプラインを示す．入力は，CBC alert rows，またはhost/process/time windowである．CLOUSEAU本来のPOI起点調査\cite{clouseau}にならい，Chief agentは初期手がかりから調査方針を生成する．Investigation agentは，確認すべき事実をQAAgent / SQL expertへの質問に変換する．QAAgent / SQL expertは，SQLite化されたエンドポイントログに問い合わせ，timestamp，source stream，PID/PPID，process tree，command line，registry path，source row idなどを返す．Final synthesisは観測結果を統合し，主行動列，補助証跡，limitationsを出力する．

\begin{figure*}[t]
\centering
\fbox{\begin{minipage}{0.92\linewidth}
\small
\centering
SOC clue / CBC alert / process-time clue\\
$\downarrow$\\
Chief agent: 調査方針を生成\\
$\downarrow$\\
Investigation agent: 確認すべき事実を質問へ変換\\
$\downarrow$\\
QAAgent / SQL expert: SQLite endpoint logsから観測行を取得\\
$\downarrow$\\
Final synthesis: 主行動列と近傍行動を分離\\
$\downarrow$\\
行動ステップ列 + 短縮行動系列 + evidence + limitations\\
$\downarrow$\\
行動主張単位の採点
\end{minipage}}
\caption{SOC調査起点から証跡付き行動列を復元するパイプライン}
\label{fig:pipeline}
\end{figure*}

本手法の設計上の要点は，仮説生成と観測ログ検索を分離することである．Chief agentは直接SQLを書かず，Investigation agentが観測済みの値に基づいて次の確認事項を作る．DBに問い合わせるのはQAAgent / SQL expertであり，Final synthesisは観測行で支えられる行動だけを主系列に含める．これにより，alert nameやreasonの言い換えを行動復元と誤認しないようにする．

\subsection{出力と評価単位}

出力は，\code{steps}，\code{sequence}，\code{evidence}，\code{nearby}，\code{limitations}からなる．\code{steps}は主体，操作，対象，コマンドライン，証跡を持つ主行動ステップである．\code{sequence}は観測された操作を短く並べた系列であり，アラート名ではなくログ上の操作を表す．\code{evidence}はsource stream，row id，timestamp，PID，command lineなどの根拠である．\code{nearby}は時刻や対象が近いが主系列に含めない行動を明示する．

採点では，gold chainを行動itemに分解する．必須itemは，ログから直接確認できる主体，操作，対象，コマンドライン，critical evidenceである．観測行だけから主張できない悪性判定，ユーザ意図，業務目的，ファイル内容の推定は必須itemに含めない．候補出力の主張は正規化され，gold itemと照合される．主指標は，行動item recall，critical evidence recall，sequence order，candidate claim precision，overclaim slot countである．

\section{実験設定}

\subsection{データとchain}

実験には，ATLASv2 benign H1 / benign-1環境をCLOUSEAUのログ探索形式へ変換したWindowsホストログを用いた．評価単位は27個のbehavior chainである．chainは，DNS packet capture batch，Python SimpleHTTPServer network，Sublime Python script execution，cmd.exe周辺操作，Discord Run Key registryの5種類のraw chain typeから構成される．分析時には，これらを明示的実行チェーン，複数段ツールチェーン，意味解釈型チェーンの3群にまとめる．

明示的実行チェーンは，親子プロセスとコマンドラインが比較的直接的に観測できるchainである．複数段ツールチェーンは，batch，packet capture tool，editor，script，file/network side effectなど，複数の道具や中間ステップを結ぶ必要があるchainである．意味解釈型チェーンは，Discord/Update.exeとRun Keyのように，アプリケーション固有またはレジストリ固有の意味解釈を必要とするchainである．

\subsection{3段階の入力条件}

表\ref{tab:stages}に，3つのstageを示す．Stage 1はCBC alert rowsを入力に含めるアラート支援条件である．Stage 2はhost，focus process，time windowのみを入力とし，DB内にはalert summary rowsを残す．Stage 3はStage 2と同じ初期入力だが，検索対象からCBC alert summary rowsを除外する．ただし，CBC EDR/NGAV event telemetry，Security，Sysmon，DNS，browser historyは残す．

\begin{table}[t]
\centering
\small
\caption{入力条件}
\label{tab:stages}
\begin{tabular}{p{0.17\linewidth}p{0.34\linewidth}p{0.37\linewidth}}
\toprule
Stage & 入力 & 意味 \\
\midrule
Stage 1 & CBC alert rows & アラート支援の復元 \\
Stage 2 & host/process/time & DB内の要約行は利用可能 \\
Stage 3 & host/process/time & alert summary rowsを検索対象から除外 \\
\bottomrule
\end{tabular}
\end{table}

この設計により，初期入力の強さと，DB内の高密度なアラート要約行の有無を分けて観察できる．Stage 2とStage 3の差は，alert summaryが最終出力として使われるかではなく，探索経路や証跡同定の足場として機能するかを調べるためのものである．

\subsection{実行と採点}

比較対象は\code{gpt-4.1-mini}と\code{gpt-5.4-mini}である．各モデルについて，27 chainを3 stageで実行し，81 runを得た．2モデル合計で162 raw runを保存し，最終採点は各モデル81 runずつに対して行った．採点対象のgoldは，起点となる観測行の特定，process actionまたはobject operationへの分解，必須item定義，critical evidence slot固定，gold chainに属さない近傍行動の除外という手順で作成した．

本稿の集計では，Stage 1/2は75 behavior steps，Stage 3は65 answerable behavior stepsを基準とする．各stepは主体，操作，対象の3つのaction itemと，1つのcritical evidence itemを持つ．Stage 3で分母が小さいのは，alert summary rowsを隠した条件で直接答えられないstepを除外し，残存telemetryから回答可能な範囲に評価を限定したためである．なお，\code{gpt-4.1-mini}の保存済み集計はlegacy形式でaction欄にcritical evidence slotを含んでいたため，本文のcross-model表ではcritical evidence分を分離した正規化action recallを用いる．

公開用パッケージには，case metadata，chain summary，gold index，runごとのモデル出力JSON，採点rubric，集計CSV/JSONを含めた．一方，元の\code{incident.db}，\code{scenario.db}，EVTX，外部データセット本体は含めない．したがって，本稿の表とrun-level挙動は検査可能であるが，完全な再実行にはローカルのATLASv2/CLOUSEAUデータベース環境が必要である．

\section{結果}

\subsection{全体結果}

表\ref{tab:overall}に，全体結果を示す．\code{gpt-5.4-mini}は81 run全体で行動item recall 83.9\%（541/645），critical evidence recall 80.9\%（174/215）を示した．一方，\code{gpt-4.1-mini}は正規化した行動item recallが47.8\%（308/645），critical evidence recallが6.0\%（13/215）であり，行動らしい説明を生成できても，必要な証跡に接続する能力が大きく不足した．

\begin{table*}[t]
\centering
\small
\caption{全体結果}
\label{tab:overall}
\begin{tabular}{lrrrrr}
\toprule
Model & Action recall & Crit. ev. & Order & Claim prec. & Overclaim \\
\midrule
\code{gpt-4.1-mini} & 47.8\% (308/645) & 6.0\% (13/215) & 18.1\% (25/138) & 34.0\% (303/892) & 589 \\
\code{gpt-5.4-mini} & 83.9\% (541/645) & 80.9\% (174/215) & 62.3\% (86/138) & 51.0\% (126/247) & 121 \\
\bottomrule
\end{tabular}
\end{table*}

注目すべき点は，\code{gpt-5.4-mini}でもcandidate claim precisionは51.0\%に留まったことである．これは，強いモデルが正解itemを多く拾える一方で，候補出力に近傍行動や重複主張も残すことを意味する．SOC支援としては広めに候補を出す価値があるが，論文上の行動列復元としては，主系列と補助証跡の分離が依然として課題である．

\subsection{Stage別結果}

表\ref{tab:bystage}にStage別結果を示す．\code{gpt-5.4-mini}はStage 1で88.9\%（200/225），Stage 2で80.4\%（181/225），Stage 3で82.1\%（160/195）の行動item recallを示した．Stage 3でも高いrecallを維持したことから，強いモデルはalert summary rowsが隠された条件でも，残存telemetryから多くのchainを復元できることが分かる．

\begin{table*}[t]
\centering
\scriptsize
\caption{Stage別結果}
\label{tab:bystage}
\begin{tabular}{llrrrrr}
\toprule
Model & Stage & Action recall & Crit. ev. & Order & Claim prec. & Overclaim \\
\midrule
\code{gpt-4.1-mini} & Stage 1 & 64.0\% (144/225) & 17.3\% (13/75) & 33.3\% (16/48) & 40.2\% (151/376) & 225 \\
\code{gpt-4.1-mini} & Stage 2 & 44.0\% (99/225) & 0.0\% (0/75) & 8.3\% (4/48) & 26.5\% (91/344) & 253 \\
\code{gpt-4.1-mini} & Stage 3 & 33.3\% (65/195) & 0.0\% (0/65) & 11.9\% (5/42) & 35.5\% (61/172) & 111 \\
\code{gpt-5.4-mini} & Stage 1 & 88.9\% (200/225) & 85.3\% (64/75) & 81.3\% (39/48) & 59.0\% (49/83) & 34 \\
\code{gpt-5.4-mini} & Stage 2 & 80.4\% (181/225) & 76.0\% (57/75) & 45.8\% (22/48) & 53.2\% (41/77) & 36 \\
\code{gpt-5.4-mini} & Stage 3 & 82.1\% (160/195) & 81.5\% (53/65) & 59.5\% (25/42) & 41.4\% (36/87) & 51 \\
\bottomrule
\end{tabular}
\end{table*}

一方で，Stage 2のorderは45.8\%であり，Stage 1の81.3\%から大きく低下した．process/time clueだけで探索を開始する場合，必要なitemに到達できても，それらを正しい順序または因果的構造として並べることは難しい．Stage 3ではclaim precisionが41.4\%へ下がり，overclaimも51へ増えた．alert summaryがない条件では，低レベルtelemetryを広く探索するため，主系列に含めるべき行動と近傍行動の境界が曖昧になりやすい．

\subsection{Scenario別結果}

表\ref{tab:scenario}に\code{gpt-5.4-mini}のscenario group別結果を示す．明示的実行チェーンではaction recall 88.2\%，critical evidence recall 86.0\%，order 79.2\%であり，親子プロセスやコマンドラインが直接観測できるchainでは高い性能を示した．複数段ツールチェーンではaction recall 81.1\%を維持したが，orderは54.8\%，claim precisionは40.3\%に低下した．複数ツールをまたぐchainでは，中間ステップの接続と過剰主張の抑制が難しい．意味解釈型チェーンは3 runのみであるが，critical evidence recall 55.6\%，order 33.3\%に留まり，registry/app固有の意味解釈が難所であることを示唆する．

\begin{table}[t]
\centering
\scriptsize
\caption{\code{gpt-5.4-mini}のscenario別結果}
\label{tab:scenario}
\begin{tabular}{lrrr}
\toprule
Group & Action & Crit. ev. & Order \\
\midrule
明示的実行 & 88.2\% & 86.0\% & 79.2\% \\
複数段ツール & 81.1\% & 78.8\% & 54.8\% \\
意味解釈 & 74.1\% & 55.6\% & 33.3\% \\
\bottomrule
\end{tabular}
\end{table}

\section{考察}

\subsection{ATLASv2上で何を評価できたか}

本研究の結果は，CLOUSEAUとATLASv2を組み合わせることで，検知後の端末行動復元を評価できることを示している．CLOUSEAUは攻撃narrative復元のための階層型調査エージェントであり，ATLAS/ATLASv2は攻撃調査文脈のデータ・手法として位置づけられてきた．本研究では，CLOUSEAUのPOI起点探索をSOCのアラート後分析に引き寄せ，ATLASv2のログ環境上で，検知ラベルの分類ではなく，アラートやprocess/time clueからログ証跡で支えられた行動chainを構成する能力を測った．この位置づけにより，攻撃ストーリ復元研究と実運用のアラートtriageの間にある，「検知後に何が起きたかを説明する」タスクを切り出せた．

また，良性chainを対象にした点も重要である．SOCでは，悪性攻撃だけでなく，管理操作，開発作業，ツール実行，アプリケーション更新なども調査対象になる．これらはアラートを発生させたり，攻撃手法に似たtelemetryを持ったりする．したがって，良性行動を証跡付きで説明できることは，誤検知低減や調査記録作成に直結する．

\subsection{アラート要約の役割}

Stage 1はアラート情報を入力に含み，Stage 2は入力から外すがDBには残し，Stage 3は検索対象からalert summary rowsを隠す．\code{gpt-5.4-mini}がStage 3でも高いrecallを維持したことは，出力がalert summaryの単純な言い換えではなく，残存telemetryからも多くの行動itemを復元できたことを示す．

ただし，Stage 3でclaim precisionが下がりoverclaimが増えたことから，alert summaryは探索の近道や境界設定として機能していた可能性がある．特に，時刻範囲，対象プロセス，対象レジストリ，検知理由が要約された行は，主系列の候補を絞る索引として働く．今後はSQL履歴を用いて，Stage 2とStage 3でどの問い合わせ段階から探索経路が分岐したかを分析する必要がある．

\subsection{残る課題}

第一の課題は，主行動列と近傍行動の分離である．強いモデルでもcandidate claim precisionは51.0\%であり，正解itemの高recallと低overclaimは同時に達成できていない．Final synthesisで，timestamp，parent-child relation，object identity，source row idの一致条件をより明示的に使い，主系列へ昇格する条件を厳しくする必要がある．

第二の課題は，順序構成である．Stage 2ではitem recallが80.4\%であるのに対し，orderは45.8\%であった．これは，必要な部品を集めることと，因果的または時系列的な行動storyへ組み上げることが別問題であることを示す．プロセス生成，子プロセスの操作，ファイル・レジストリ・ネットワーク副作用を，単なる列挙ではなく制約付き系列として扱う必要がある．

第三の課題は，意味解釈型chainである．Discord Run Keyのような事例では，レジストリpath，値名，データ，アプリケーション更新機構の関係を解釈する必要がある．これは低レベルログの検索だけではなく，対象オブジェクトの意味を誤らずに説明する能力を要求する．

本評価の限界として，対象は単一のATLASv2 benign環境から構成した27 chainであり，組織横断のSOCログや攻撃キャンペーン全体を代表するものではない．また，Stage 3はalert summary rowsのみを検索対象から除外する条件であり，EDR telemetry全体を除外した条件ではない．このため，本稿の結論は「alert summaryなしでも一定の行動復元が可能」という範囲に留め，「任意の低レベルログだけで十分」とは主張しない．

\section{おわりに}

本稿では，CLOUSEAUを基盤とするエージェント型ログ探索手法により，SOC調査起点からWindowsエンドポイントログを探索し，証跡付き行動列を復元する課題を，ATLASv2 benign H1 / benign-1由来の27 behavior chainで評価した．\code{gpt-5.4-mini}は81 run全体で行動item recall 83.9\%，critical evidence recall 80.9\%を示し，alert summary rowsを隠したStage 3でも82.1\%の行動item recallを維持した．一方で，candidate claim precision，overclaim，sequence orderには課題が残った．

本研究の立ち位置は，CLOUSEAUの攻撃narrative復元をそのまま再評価するものでも，ATLAS/ATLASv2を用いた攻撃検知評価でもなく，検知後のSOCインシデント分析における行動列復元評価である．今後は，SQL探索履歴の分析，主系列昇格条件の強化，順序制約を持つsynthesis，意味解釈型chainの拡充により，調査者がそのまま検証できる行動storyの品質を高める．

\begin{thebibliography}{9}
\bibitem{clouseau}
A. Aldaihan, F. Alotaibi, and S. Maffeis,
``Clouseau: A Hierarchical Multi-Agent Approach For Autonomous Attack Investigation,''
2025. \url{https://www.doc.ic.ac.uk/~maffeis/papers/acsac25.pdf}

\bibitem{sysmon}
Microsoft,
``Sysmon - Sysinternals,''
\url{https://learn.microsoft.com/en-us/sysinternals/downloads/sysmon}
（参照 2026-06-09）．

\bibitem{react}
S. Yao, J. Zhao, D. Yu, N. Du, I. Shafran, K. Narasimhan, and Y. Cao,
``ReAct: Synergizing Reasoning and Acting in Language Models,''
arXiv:2210.03629, 2022.

\bibitem{deeplog}
M. Du, F. Li, G. Zheng, and V. Srikumar,
``DeepLog: Anomaly Detection and Diagnosis from System Logs through Deep Learning,''
Proc. ACM CCS, pp. 1285--1298, 2017.

\bibitem{atlas}
A. Alsaheel, Y. Nan, S. Ma, L. Yu, G. Walkup, Z. B. Celik, X. Zhang, and D. Xu,
``ATLAS: A Sequence-based Learning Approach for Attack Investigation,''
Proc. USENIX Security, pp. 3005--3022, 2021.

\bibitem{atlasv2}
A. Riddle, K. Westfall, and A. Bates,
``ATLASv2: ATLAS Attack Engagements, Version 2,''
arXiv:2401.01341, 2023.
\end{thebibliography}

\end{CJK}
\end{document}
